\chapter*{\large ВВЕДЕНИЕ}
\addcontentsline{toc}{chapter}{ВВЕДЕНИЕ}

Разрабатывалось авторское приложение \textbf{Lyrata}~\appa{2}: пользователь задаёт свой текст и набор аудио (плейлист); во время работы над ним возникла необходимость в средстве автоматической подсказки, какой именно трек из этого набора разумнее сопоставить с текущей позицией чтения, без ручной расстановки по каждому фрагменту и без выбора треков полностью наугад --- последнее обычно плохо согласуется с восприятиями текста.

В интерфейсе удалось сделать извлечение позиции чтения с точностью до нескольких слов, что позволяет динамически менять громкость или переключать аудио при изменении участка текста на экране. Для задачи данной работы необходима модель машинного обучения~\cite{nikolenko2018}, восстанавливающая описание звука по короткому текстовому фрагменту вокруг текущей точки при ограничении кандидатов только аудио дорожками пользовательского плейлиста.

Обучающие данные для модели получены из \textit{бесплатных игр с открытым исходным кодом}, разрабатываемых на движке Ren'Py~\cite{srcrenpydocs}. Такие игры относятся к жанру \textit{визуальных новелл}: пользователь последовательно читает диалоги и короткие описания, а автор заранее прописывает, какая музыка или звуковая дорожка должна звучать в каждый момент. Это даёт возможность автоматически извлекать пары ``фрагмент текста --- связанный трек'' по файлам сценария и по-прежнему вручную подчищать или дополнять разметку там, где автоматика ошибается.

\textit{Объект исследования} --- сопоставление текстовых эпизодов игр на Ren'Py и числовых описаний музыкальных фрагментов.

\textit{Предмет исследования} --- методы обучения нейросетевых регрессионных моделей от русскоязычных текстовых эмбеддингов к вектору музыкальных признаков и оценка качества подстановки трека в списке кандидатов одной игры.

\textit{Цель} --- исследовать, возможно ли своими силами построить модель, которая по тексту, который передаёт сам пользователь (фрагмент вокруг текущей позиции чтения), помогает подбирать треки из его собственной подборки к этому месту в тексте --- сужая или упорядочивая кандидатов, а не заменяя выбор человека.

\textit{Методы} --- работа разделена на взаимосвязанные этапы: сбор и подготовка данных для обучения; построение регрессионной модели и её обучение с учителем в PyTorch; сравнение нескольких архитектур (линейная модель, MLP, MLP с лагами и авторегрессией, LSTM, GRU, TCN, трансформер по окну~\cite{vaswani2017}) на едином протоколе с отложенными играми~\cite{nikolenko2018}. Оценка качества --- по метрикам регрессии и по подстановке трека в списке кандидатов; численные результаты сняты с сохранённых прогонов в дереве \path{rubert_runs/} репозитория:~\appa{5},~\appa{1}.

\textit{Новизна} работы для курсового уровня --- в связке прикладной задачи чтения под музыку, русскоязычных эмбеддингов RuBERT (адаптация архитектуры BERT~\cite{devlin2018}, см.~\cite{kuratov2019}) и высокоразмерного целевого описания аудио~\cite{librosa} с проверкой не только ошибки регрессии, но и качества ранжирования треков внутри одной игры~\cite{nikolenko2018}.

Глава~\ref{ch:theory} задаёт сценарий использования, объясняет устройство данных Ren'Py и форматы представления текста и музыки, а также обзор использованных архитектур. Глава~\ref{ch:dataset} посвящена сбору датасета. Глава~\ref{ch:training} содержит постановку обучения, определения метрик, протокол экспериментов и численное сравнение моделей.
